






\section{Introduction}
\label{sec:intro}
Recent retrieval-augmented image captioning models have shown success in strong image captioning performance while reducing model parameters by retrieving related captions for a given image~\citep{ramos2023smallcap,rag-transformer,yang2023re}. These models use retrieved information as additional context besides the input image.
However, similar to retrieval-augmented language models \citep{ret-robust23}, image captioning models enhanced with retrieval can sometimes be misled by irrelevant information.
For example, in Figure~\ref{fig:ex} the captioning model is misled by the token ``elephant'' in the retrieved captions, and generates captions that do not match the given image.

For retrieval-augmented language models, \citet{ret-robust23} have studied the cases where retrieval misled the model prediction, and address this problem with a retrieval-robust LLM by continuous training with synthetic data for question answering tasks. However, in their approach, the retrieval system returns only one passage at each step. Considering that LLMs can be sensitive to the order of prompts~\cite{lu-etal-2022-fantastically}, the robustness of using multiple retrieved results has not been fully studied. Evaluating and improving the robustness of retrieval-augmented image captioning models remains under-explored, specifically when the model is augmented with multiple retrieved results.




    \begin{figure}[t]
        \includegraphics[width=\columnwidth]{figs/fig1_comparison.pdf}
        \caption{Comparison of generated image captions that are predicted without retrieval, misled by retrieval, and predicted with a more retrieval-robust model. The retrieval-augmented model generates the token ``elephant'', which appears in 3/4 of the retrieved captions.\label{fig:ex}}
    \end{figure}

To bridge this gap in the literature, we study the robustness of the \textsc{SmallCap} retrieval-augmented captioning model~\citep{ramos2023smallcap}. By the definition of retrieval robustness proposed in~\citet{ret-robust23}, retrieved context should boost model performance when relevant, and should not adversely affect it when irrelevant. We thoroughly examine the robustness of
the model with regards to the order of the retrieved captions, and the relevance of the retrieved content. We also present a novel analysis of model behaviour based on \textit{majority voting}, supported by input attribution and attention analyses to investigate how the retrieved tokens influence the model generation. And finally, inspired by~\citet{hoang2022improving}, we propose to sample the retrieved captions from a larger list during training to prevent the model from overfitting to the top relevant captions. Our evaluation shows improved model robustness and better out-of-domain generalization. 

The main findings of this paper are: \textbf{1)} We study the robustness of an existing retrieval-augmented captioning model \textsc{SmallCap} and find it is not robust to processing randomly retrieved content. \textbf{2)} We identify that tokens that frequently occur in the retrieved captions, i.e. majority tokens, have high attribution scores with regard to the tokens generated by the model. This phenomenon suggests heightened sensitivity and copying. \textbf{3)} Training with sampled retrieved captions from a larger list instead of with fixed top-k relevant captions improves model robustness, yielding better generalization and out-of-domain performance.\footnote{We release the code at \url{https://github.com/lyan62/RobustCap}} 




    

